\section{OpenAI、Navier--Stokes問題への解を発表}
OpenAIは、未公開の次世代モデルと約1万の協調エージェントがNavier--Stokesの存在と滑らかさ問題に対する解析的証明とLean形式化を生成したと発表した。公式説明では88時間で解へ到達し、追加17時間でAstraを用いたLean検証を行ったとしている。一方、数学的正しさやClay賞の認定はまだ確定しておらず、報道も「OpenAIが解いたと主張」と慎重に扱っている。
\dailyxsource{OpenAI (@OpenAI)}{https://x.com/OpenAI/status/2097374640582668336}
\section{Meta、個人エージェントMuseをローンチ}
MetaはMuse Spark 1.3を基盤とする個人エージェントMuseを発表した。公式投稿では、生活全般のタスク代行、アプリ接続のユーザー制御、購入やメール操作時の確認、Sentinelエージェント、資格情報の分離保管などを説明している。無料枠は週1億トークンまでとされるが、対応地域の完全な一覧はGrok収集では確認されていない。
\dailyxsource{Muse (@Muse)}{https://x.com/Muse/status/2097399178376671666}
\section{ChatGPT Images 2.5とFlare／Sunburst}
OpenAIはChatGPT Images 2.5をChatGPT、ChatGPT Work、Codexのデスクトップ・モバイル・Webへ全ユーザー向けに展開すると発表した。APIにはGPT-Image-2.5 FlareとSunburstを追加し、Flareは速度と編集、Sunburstはより精細な生成を重視すると説明している。専用ブログURLや価格はGrok収集では未確認である。
\dailyxsource{OpenAI (@OpenAI)}{https://x.com/OpenAI/status/2097394956457623964}
\section{Astra、CodexとChatGPT Workへ全面展開}
OpenAIはGPT-6 AstraがPlus、Pro、Business、Enterpriseユーザー向けにCodexとChatGPT Workでfully rolled outになったと投稿した。前日までの段階的提供から、対象有料プランでの本格展開へ進んだ形である。今回のDaily版では製品ドキュメントの追加確認は行わず、公式投稿の範囲を記録する。
\dailyxsource{OpenAI (@OpenAI)}{https://x.com/OpenAI/status/2097431322117476423}
\section{Inception、Mercury 2.5を発表}
Inceptionは拡散LLM Mercury 2.5を発表した。公式ブログはMercury 2比でintelligence 40\%向上、汎用NVIDIA GPU上で1,107 tok/s、260K context、入力0.20ドル／出力0.75ドル per million tokensと説明し、ローンチ時80\%オフも掲げている。API、OpenRouter、Basetenで利用可能と公式Xが述べた。
\dailyxsource{Inception (@\_inception\_ai)}{https://x.com/_inception_ai/status/2097365772289151417}
\section{Cohere、decode megakernelをOSS公開}
CohereはLLM decode全体を単一カーネルに融合するmegakernelベースのサービング実装を公開した。公式投稿ではNorth Mini CodeをBF16・1台のH100で動かした条件で、vLLM比最大1.58倍などの速度向上を主張している。ブログとGitHubリポジトリが公開されているが、他モデルや他GPUでの再現は今回のDaily工程では扱わない。
\dailyxsource{Cohere (@cohere)}{https://x.com/cohere/status/2097410772355666393}
\section{Cursor、Muse Spark 1.3を提供開始}
CursorはMetaのMuse Spark 1.3をIDE内で利用可能にしたと発表した。長時間のコーディングエージェント用途と価格性能を訴求し、CursorBenchへのリンクも示している。Meta側のMuse発表とほぼ同時期に、個人エージェントと開発者向けIDEの双方へSpark 1.3が広がった。
\dailyxsource{Cursor (@cursor\_ai)}{https://x.com/cursor_ai/status/2097402609531236708}
\section{Sam Altman、GPT-6関連イベントを告知}
Sam Altmanは9月16日にサンフランシスコでGPT-6について語り交流するイベントを開催すると告知した。応募締切は9月10日で、GPT-5.5時にも同様の集まりを行ったと述べている。新モデルの発表そのものではなく、ローンチ後のコミュニティイベントとして観測された話題である。
\dailyxsource{Sam Altman (@sama)}{https://x.com/sama/status/2097404861642137851}
\section{Anthropicの抗生物質発見説が拡散}
Anthropicが耐性菌に有効な新クラス抗生物質を発見し、IPO時に発表する予定だという噂がXで拡散した。元投稿自身がRumourと明記しており、観測窓内にAnthropic公式アカウントによる確認は見つかっていない。Daily版では企業・研究上の確定情報としてではなく、未検証の噂が流通した事実として扱う。
\dailyxsource{billy (@billyhumblebrag)}{https://x.com/billyhumblebrag/status/2097397048559440065}
\section{Tamarindの分子AIモデルルータ}
SynBioBetaは、TamarindのMolecular AI Model Routerを紹介した。400超のモデル基盤上で入力ごとに適したモデルを選ぶ構想で、第一弾は分子構造予測向けとされる。Tamarind公式Xの一次投稿はGrok収集では確認されておらず、専門メディア投稿とブログURLを基にした話題として記録する。
\dailyxsource{SynBioBeta (@SynBioBeta)}{https://x.com/SynBioBeta/status/2097433153204097354}
\section{クロスモデルKV cache変換の論文解説}
NVIDIA系の最近の論文として、あるモデルのKV cacheを同一ファミリーの別モデル向け表現へ訓練なしの線形写像で変換する手法が解説された。投稿では再プリフィルより2.7--25倍高速で、6ペア中4ペアで単独精度の73--98\%を維持したとされる。一方、正式なarXiv URLや発表日時は未確認で、技術解説投稿として扱う。
\dailyxsource{Akshay (@akshay\_pachaar)}{https://x.com/akshay_pachaar/status/2097421509220561028}
\section{Astraに計算予算を与えた個人実験}
Harry GaoはAstraにModal上で200ドル相当の計算予算を与え、過去の拡散モデル論文を改善させる実験の途中経過を報告した。21時間、28 GPU-hours、117.04ドルを使い、後続ベンチで誤差19.2\%減と自己報告している。再現物やPRは未提示であり、詳細writeup待ちの個人実験として記録する。
\dailyxsource{Harry Gao (@hrygao)}{https://x.com/hrygao/status/2097112684441546811}
